\subsection*{S4. Explicit \texttt{null} and absent fields}

JSON distinguishes three states where OCaml's \texttt{option} offers two: a
key present holding a value, a key present holding \texttt{null}, and a key
absent altogether.

\begin{lstlisting}[style=json]
{ "a": 1, "b": 2 }        -- b present with a value
{ "a": 1, "b": null }     -- b present, explicitly null
{ "a": 1 }                -- b absent
\end{lstlisting}

Two independent facts are observable per path: whether the key is ever absent,
and whether the value is ever \texttt{null}. Preserving both would require a
three-state type; collapsing them fits \texttt{option} directly.

\options
\begin{enumerate}[label=\textbf{(\alph*)}]
  \item \textbf{Collapse.} \texttt{null} and absence both mean ``no value''.
        A path that is ever absent or ever null becomes \texttt{option}, and
        \texttt{None} covers both. \selected
  \item \textbf{Keep them apart} in the type, as \texttt{int option option} or
        as a generated three-state variant
        \texttt{type 'a field = Absent | Null | Value of 'a}. Exact, but
        heavier to consume and it propagates into the Phase~2 effect
        signatures.
  \item \textbf{Distinguish in the schema only.} Record both facts during
        inference but generate a single \texttt{option}, so the distinction
        survives in the schema and the documentation but not in the type.
  \item \textbf{Reject explicit \texttt{null}.} Only absence means ``no
        value''; a corpus containing \texttt{null} is an error.
\end{enumerate}

\decision{Collapse. A path whose key is absent in some document, or whose
value is \texttt{null} in some document, becomes \texttt{option}; a path
always present and never null carries no \texttt{option} at all.}

\begin{lstlisting}[style=json]
{ "a": 1, "b": 2 }
{ "a": 1, "b": null }
{ "a": 1 }

  -> a : int              (always present, never null)
  -> b : int option       (None covers both null and absent)
\end{lstlisting}

A path needs only one document to make it partial; the rest may all supply
values:

\begin{lstlisting}[style=json]
{ "x": 1 }   { "x": 2 }   { "x": 3 }   { }

  -> x : int option       (absent once, so partial everywhere)
\end{lstlisting}

\rationale{This is the two-state framing the design note already uses when it
says fields are classified as total or partial, and it is what comparable
tools do. It keeps the generated modules directly consumable, and it keeps the
Phase~2 effect signatures as simple as the underlying scalar allows.}

\limitation{The two absent-value cases become indistinguishable once shredded.
A document carrying \texttt{"b": null} and a document omitting \texttt{b}
produce identical rows, so the original cannot be recovered exactly: any
round-trip property will hold only up to this collapse. A related consequence
is that reconstruction must pick a canonical form when it emits a
\texttt{None} --- either \texttt{"b": null} or omitting the key --- and that
choice is not made here.}
